% !TeX program = xelatex
\documentclass[10pt]{article}

\usepackage{hyperref}
\usepackage{xcolor}
\usepackage{calc}
\usepackage{graphicx}
\usepackage{tikz}
\usepackage{fontspec}
\usepackage{fontawesome5}
\usepackage{titlesec}
\usepackage{enumitem}
\usepackage{fancybox}
\usepackage{xeCJK}

\hypersetup{hidelinks}
\urlstyle{same}

%%%%%%%%%%%%%%%%%%%%
% 设置
%%%%%%%%%%%%%%%%%%%%

\providecommand{\ResumeItemSep}{0.12em}
\providecommand{\ResumeTopSep}{0.18em}
\providecommand{\ResumeArrayStretch}{1.12}
\providecommand{\ResumeLineSpread}{1.08}
\providecommand{\ResumeSectionBefore}{0.55em}
\providecommand{\ResumeSectionAfter}{0.28em}

\setlength{\parindent}{0pt}
\setlength{\parskip}{0pt}
\pagenumbering{gobble}
\setlist[itemize]{
    leftmargin=1.35em,
    itemsep=\ResumeItemSep,
    topsep=\ResumeTopSep,
    parsep=0pt,
    partopsep=0pt
}
\setlist[enumerate]{leftmargin=*}
\renewcommand{\arraystretch}{\ResumeArrayStretch}
\linespread{\ResumeLineSpread}

\titleformat{\section}
  {\Large\bfseries\raggedright}
  {}{0em}
  {}
  [{\color{secondarycolor}\titlerule}]
\titlespacing*{\section}{0pt}{\ResumeSectionBefore}{\ResumeSectionAfter}

\usepackage[
    a4paper,
    left=1.1cm,
    right=1.1cm,
    top=0.75cm,
    bottom=0.85cm,
    nohead
]{geometry}

% 字体设置
% 说明：Noto Serif SC 没有原生斜体字型。将 CJK 的 italic 字型映射到常规字形
%（与过去“找不到斜体、默认替换为常规字形”的实际渲染一致），可消除
% “Font shape `TU/NotoSerifSC(0)/m/it' undefined” 告警；如需真正斜体观感，
% 可在此行追加 ItalicFeatures={FakeSlant=0.18}。
\setCJKmainfont[
    Path=fonts/,
    Extension=.otf,
    BoldFont=*-Bold,
    ItalicFont={NotoSerifSC},
]{NotoSerifSC}

% 自定义颜色
\definecolor{primarycolor}{RGB}{200, 60, 60}
\definecolor{secondarycolor}{RGB}{200, 68, 28}

\newlength{\iconwidth}
\setlength{\iconwidth}{1.5em}

\newcommand{\sectionicon}[2]{\makebox[\iconwidth][c]{\color{primarycolor}{#1}}\quad #2}

\newcommand{\projectheading}[3]{%
    \noindent\makebox[\textwidth][s]{%
        \makebox[0pt][l]{{\large \textbf{#1}}}%
        \hfill
        \makebox[0pt][c]{{\normalsize \textbf{#2}}}%
        \hfill
        \makebox[0pt][r]{#3}%
    }%
}

% 通用个人信息：修改这里即可同步到所有岗位模板。
\newcommand{\ResumeName}{无绝}
\newcommand{\ResumeBirthDate}{2002.10}
\newcommand{\ResumeHometown}{湖北咸宁}
\newcommand{\ResumeSchool}{重庆邮电大学}
\newcommand{\ResumeEnglishLevel}{CET-6}
\newcommand{\ResumeEmail}{250xxxxxxx@qq.com}
\newcommand{\ResumePhone}{198xxxxxxxx}
\newcommand{\ResumeHomepageLabel}{github.com/Dithob}
\newcommand{\ResumeHomepageUrl}{https://github.com/Dithob}
\newcommand{\ResumePhoto}{images/xjpic.jpg}
\newcommand{\ResumeHeaderImage}{images/foot.png}

% 通用教育背景：修改这里即可同步到所有岗位模板。
\newcommand{\ResumeGraduateSchool}{重庆邮电大学}
\newcommand{\ResumeGraduateMajor}{计算机技术}
\newcommand{\ResumeGraduateDegree}{硕士}
\newcommand{\ResumeGraduateDate}{2024.09--2027.06}
\newcommand{\ResumeGraduateAwards}{特等学业奖学金 $\times$ 1、二等学业奖学金 $\times$ 1}

\newcommand{\ResumeUndergraduateSchool}{重庆邮电大学}
\newcommand{\ResumeUndergraduateMajor}{计算机科学与技术}
\newcommand{\ResumeUndergraduateDegree}{本科}
\newcommand{\ResumeUndergraduateDate}{2020.09--2024.07}
\newcommand{\ResumeUndergraduateAwards}{三等学业奖学金；GPA: 3.28/4.0（专业前 25\%）}

\newcommand{\ResumeEducationBlock}{%
    {\textbf{\ResumeGraduateSchool}}，\ResumeGraduateMajor，\textit{\ResumeGraduateDegree} \hfill \ResumeGraduateDate
    \begin{itemize}
        \item \textbf{荣誉奖项}：\ResumeGraduateAwards
    \end{itemize}

    \vspace{0.2em}
    {\textbf{\ResumeUndergraduateSchool}}，\ResumeUndergraduateMajor，\textit{\ResumeUndergraduateDegree} \hfill \ResumeUndergraduateDate
    \begin{itemize}
        \item \textbf{荣誉奖项}：\ResumeUndergraduateAwards
    \end{itemize}
}


%%%%%%%%%%%%%%%%%%%%
% 文章内容
%%%%%%%%%%%%%%%%%%%%

\newcommand{\ResumeTargetRole}{AI 应用开发（Agent 质量效能）}
\providecommand{\ResumeTargetRole}{求职方向待定}

\newcommand{\ResumeContact}{%
    \footnotesize
    \textcolor{white}{%
        \faEnvelope \quad 邮箱：\href{mailto:\ResumeEmail}{\ResumeEmail}
        \hspace{3.2em}
        \faPhone \quad 电话：\ResumePhone
        \hspace{3.2em}
        \faGithub \quad 主页：\href{\ResumeHomepageUrl}{\ResumeHomepageLabel}
    }%
}

\newcommand{\ResumeContactBar}{%
    \begin{tikzpicture}[remember picture, overlay]
        \node[anchor=north, inner sep=0pt](headercontacts) at (current page.north){%
            \includegraphics[width=\paperwidth]{\ResumeHeaderImage}
        };
        \node[anchor=center] at(headercontacts.center){\ResumeContact};
    \end{tikzpicture}
    \vspace{-1.15em}
}

\newcommand{\ResumeProfileSection}{%
    \section[个人信息]{\sectionicon{\faAddressCard}{个人信息}}
    \begin{tabular*}{\textwidth}{@{\extracolsep{\fill}}lll}
        \textbf{姓\qquad 名}：\ResumeName & \textbf{出生年月}：\ResumeBirthDate & \textbf{籍\qquad 贯}：\ResumeHometown \\
        \textbf{学\qquad 校}：\ResumeSchool & \textbf{英语水平}：\ResumeEnglishLevel & \textbf{求职方向}：\ResumeTargetRole
    \end{tabular*}
}

\newcommand{\ResumeEducationSection}{%
    \section[教育背景]{\sectionicon{\faGraduationCap}{教育背景}}
    \ResumeEducationBlock
}

\newcommand{\ResumeProfileAndEducation}{%
    %%%%%%%%%%%%%%%%%%%%
    % 个人信息与教育背景
    %%%%%%%%%%%%%%%%%%%%

    \begin{minipage}[t]{0.79\textwidth}
        \ResumeProfileSection

        \ResumeEducationSection
    \end{minipage}
    \hfill
    \begin{minipage}[t]{0.14\textwidth}
        \vspace{-0.35em}
        \setlength{\fboxsep}{0pt}
        % 双线框 (\doublebox) 会给照片额外加约 4pt 边框宽；图片若占满 \linewidth 会触发
        % “Overfull \hbox (3.99988pt too wide)” 告警，故图片宽度预留 4pt，使含框总宽恰好等于栏宽。
        \doublebox{\includegraphics[width=\dimexpr\linewidth-4pt\relax]{\ResumePhoto}}
    \end{minipage}
}


\begin{document}

    % 顶部联系人栏
    \ResumeContactBar
    \ResumeProfileAndEducation
    % \vspace{0.45em}

    %%%%%%%%%%%%%%%%%%%%
    % 专业技能
    %%%%%%%%%%%%%%%%%%%%

    \section[专业技能]{\makebox[\iconwidth][c]{\color{primarycolor}{\faWrench}}\quad 专业技能}
    \begin{itemize}
        \item \textbf{AI 应用开发（Agent 赋能）}：
        熟悉 Dify、LangChain 等 LLM 应用开发框架，能够围绕业务场景设计 RAG 问答、工具调用、多 Agent 协作与结构化输出流程，并使用 FastAPI 封装大模型服务接口；
        理解 Harness Engineering、Agent Memory、MCP 协议、Skill/Subagent 编排等 Agent 核心机制，专注以 AI/Harness/Agent 为测试、QA 与需求推进赋能，构建“需求理解 $\rightarrow$ 用例生成 $\rightarrow$ 自动执行 $\rightarrow$ 失败自愈”的质量效能闭环；熟练应用 Claude Code、Codex、CodeBuddy 等 AI 编程工具。
        % \item \textbf{RAG 与检索增强}：熟悉文档 ETL、Chunk 切分、Embedding、Milvus 向量检索、BM25 关键词召回、RRF 融合与 Rerank 精排；能够构建知识库检索链路，并通过相似度阈值、拒答策略和 Bad Case 回流降低幻觉风险。
        \item \textbf{模型微调与推理优化}：
        理解 Transformer、自注意力 等底层架构原理；
        % 理解 Transformer、自注意力 等底层架构原理、掌握 Prompt Engineering、CoT 等提示工程技巧；
        了解 SFT 与 PPO 等训练范式，熟悉 LoRA 等微调方法，具备 Qwen 微调与模型 API 集成经验，熟悉 vLLM 推理部署以及KV Cache 管理与连续批处理调优。
        \item \textbf{计算机基础与工程能力}：
        掌握数据结构、计算机网络、数据库原理等基础知识；
        掌握 Python 语言；熟悉 MySQL/SQLite 数据库设计与优化，理解 Redis 核心数据结构及缓存架构；
        熟悉Linux常用操作，掌握 Docker容器化部署与 Git 版本控制；
        了解测试理论，能熟练运用 Postman/JMeter进行接口联调与性能压测，具备完整的环境隔离与工程交付能力。
        % \item \textbf{全栈与 AI 协作开发}：
        % 熟悉 \textbf{Vue.js/Element} 前端生态与 \textbf{Spring Boot/MyBatis} 后端框架，具备前后端分离项目的独立闭环开发能力；
        % 熟练应用 Claude Code、Codex、Qoder 等 AI 编程工具，并在项目中实践 \textbf{MCP Server、Agent SKILL} 等前沿范式，大幅提升开发效能。
    \end{itemize}

    % \vspace{0.45em}

    %%%%%%%%%%%%%%%%%%%%
    % 实习经历
    %%%%%%%%%%%%%%%%%%%%

    \section[实习经历]{\makebox[\iconwidth][c]{\color{primarycolor}{\faBriefcase}}\quad 实习经历}

    \projectheading{腾讯地图}{AI 测试开发实习生（Agent 工程方向）}{2026.07--2026.10}
    \begin{itemize}
        \item \textbf{核心职责}：参与地图 AI Native 质量效能体系建设，负责 Agent 编排设计、Skill 工作流与 MCP 工具落地及业务域用例生成，以 AI/Harness/Agent 为测试提效与需求推进赋能，打通“需求理解 $\rightarrow$ 代码生成 $\rightarrow$ 审查修复 $\rightarrow$ 精准回归”全流程的工程化接入。
        \item \textbf{Agent 设计与可靠性护栏}：参与 13 个 Skill 五层编排体系与 3 个自研 MCP 服务设计，落地 Given/When/Then 意图锁定、探索门禁、单步探针验证、九维度自审与探索预算人工接管等护栏机制，保障 LLM 规模化生成代码的质量与可控性。
        \item \textbf{多 Agent 交叉检查与 Skill 评测}：使用多 Agent 交叉检查流程完成需求用例生成（多智能体独立生成 $\rightarrow$ 交叉比对合并，产出 85 条 P0/P1/P2 分级用例），设计人工复核口径标注 39 处待确认项，形成“AI 生成 + 人工质量判断”闭环；为地图数据查询、地图渲染两类 Skill 规划 120 条评测用例，覆盖能力调用、参数边界、异常输入与安全约束，将 Agent/Skill 能力验证结构化。
    \end{itemize}

    \vspace{0.25em}

    \projectheading{重庆啄木鸟网络科技有限公司}{算法工程师}{2025.06--2025.09}
    \begin{itemize}
        \item \textbf{核心职责}：参与家电售后场景的大模型应用研发，覆盖需求分析、数据清洗、模型调研与训练、Prompt/RAG 流程优化、接口设计与测试、FastAPI 服务封装及部署联调等环节。
        \item \textbf{小啄 GPT 智能回复模型}：
        针对多轮客服对话中的上下文丢失、槽位遗漏与知识检索不准问题，设计“短期记忆 + 长期摘要 + 结构化状态”机制，用 JSON 维护家电类型、故障现象、预约时间等关键状态；
        针对生成环节，对 Qwen 模型进行微调，注入家电维修领域知识与标准客服话术，有效抑制模型幻觉，使最终回复的业务采纳率提升约 10\%；
        结合 Dense/BM25 双路召回、RRF 融合与 Rerank 精排构建维修知识检索链路，使长尾 Query 检索 Hit Rate 提升约 20\%，提升长轮次对话的流程控制与回复可靠性。
        % \item \textbf{智能产品识别与意图理解}：
        % 参与智能产品识别、用户意图识别等模块迭代，设计工作流从用户咨询中抽取品牌、型号、品类、故障现象等结构化槽位，并判断对话意图类型；
        % 通过 问题重写 规范口语化/残缺 问询，结合语义匹配、RAG 知识库检索与动态追问策略完成候选产品召回、槽位补全和低置信度确认；
        % 支撑 1000+ 类产品识别和 20+ 类用户意图实时识别，为后续精准问答、工单分发与人工兜底提供结构化依据。
    \end{itemize}

    % \vspace{0.45em}

    %%%%%%%%%%%%%%%%%%%%
    % 项目经历
    %%%%%%%%%%%%%%%%%%%%

    \section[项目经历]{\makebox[\iconwidth][c]{\color{primarycolor}{\faProjectDiagram}}\quad 项目经历}
    % AI 应用开发 / 大模型应用岗
    \projectheading{LLM Agent 质量效能工程（腾讯地图）}{核心参与（Agent 编排落地）}{2026.07--2026.10}
    \begin{itemize}
        \item \textbf{项目概述}：基于 CodeBuddy（Skill + Subagent + MCP + Rules + Memory）构建 LLM Agent 质量效能体系，让大模型稳定完成“自然语言需求 $\rightarrow$ 可执行代码生成 $\rightarrow$ 真机执行 $\rightarrow$ 失败自愈 $\rightarrow$ 知识沉淀”的闭环，替代传统人工编写脚本模式，为测试与需求推进提效；覆盖 17 个业务域、600+ 生成代码单元。
        \item \textbf{多 Agent 编排与可靠性设计}：按职责五层编排 13 个 Skill 与 3 个自研 MCP 服务；以 Given/When/Then 意图锁定防止模型为“跑通”而偷偷降标，探索门禁 + 单步验证 + 九维度自审控制生成质量，设置探索预算与人工接管机制实现人机协同止损，单任务生成约 10--15 分钟（人工需 1--2 小时）。
        \item \textbf{RAG 与知识工程}：设计纯文件知识图谱（4 类节点 + 6 类边，按域加载降噪）与同源代码库静态推理（23 个仓库 UI 索引快照），让 Agent 生成前先检索理解研发代码；以 AST 三道闸扫描 1500+ 函数，结合 jieba + BM25 + 120+ 受控能力词表实现公共能力语义检索与自动沉淀，解决 LLM 规模化生成时“重复造轮子”问题。
        \item \textbf{多模态感知与自愈闭环}：融合 UIAutomator XML / RapidOCR / OmniParser(YOLO) / OpenCV 模板四层感知，链式降级 + 命中率自适应排序解决无 resource-id 控件定位问题；失败后基于证据包（截图/trace/标注图）自动再生成补丁，形成 Self-heal 闭环；scrcpy 持久流将截图耗时由约 1s 降至 30--100ms。
    \end{itemize}

    \vspace{0.35em}
    \projectheading{基于 RAG 的医药数据分析问答助手}{项目主要成员}{2025.10--2026.02}
    % 算法 / NLP / LLM 工程岗
    % \projectheading{医药供应链 Text-to-SQL 数据分析助手（校企合作）}{项目主要成员}{2025.10--2026.02}
    \begin{itemize}
        \item \textbf{项目目标}：面向药品采购、零售、库存等全链路数据及 800 余张业务表，基于 Dify、LangChain、Milvus 与 MySQL 建设 LLM 数据分析问答助手，将销售人员的自然语言问题转化为可执行 SQL 与可读分析结论。
        \item \textbf{应用架构设计}：负责 LLM 应用核心链路设计，基于 Dify 完成 工作流 编排与模型接入，设计“问题理解--Schema 检索--SQL 生成--结果解释--人工反馈”业务链路，降低业务人员直接理解复杂数据库结构的门槛。
        \item \textbf{表结构理解与字段召回}：对数据字典进行元数据增强，将表名、字段名、字段含义和常见枚举值拼接为富文本后向量化；结合 业务规则过滤、问题类型判断、样例对与 CoT ，引导模型优先选择相关表和字段，提高 SQL 生成稳定性。
        % \item \textbf{表结构理解与字段召回}：对数据字典进行元数据增强，将表名、字段名、字段含义和常见枚举值拼接为富文本后向量化；结合 HNSW 索引、标量过滤、规则化问题分类、样例对与 CoT ，引导模型优先选择相关表和字段，提高多表关联、专业术语和时间/分类条件下的 SQL 生成稳定性。
        \item \textbf{评估与自修正闭环}：引入 SQL 预执行与 自纠错 机制，捕获 MySQL Error 信息后回传给 LLM 重写；构建约 200 条 Golden Dataset，结合 RAGAS 指标、执行准确率 与人工反馈 Bad Case 回流评估检索和 SQL 结果，执行成功率由约 70\% 提升至 90\%+，项目获重庆市人工智能大赛市级银奖。
    \end{itemize}

    \vspace{0.35em}
    \projectheading{多模态家电识别与知识库入库系统}{项目主要负责人}{2025.06--2025.09}
    % \projectheading{基于多模态识别的家电知识库入库系统}{项目主要负责人}{2025.06--2025.09}
    \begin{itemize}
        \item \textbf{项目目标}：面向家电售后知识库建设场景，针对图片/OCR 识别信息不完整、长尾产品文本特征稀疏和重复入库问题，设计“识别--补全--校验--去重--入库”的自动化流程，降低人工录入和知识维护成本。
        \item \textbf{识别与字段补全}：使用 Docker/Miniconda 统一部署 OCR/视觉大模型环境，基于 Dify 编排“图片识别--联网搜索--字段抽取--置信度校验”多 Agent 流程；结合 Query Rewrite 和动态 Prompt 模板补全品牌、型号、规格、品类等产品字段。
        % \item \textbf{检索增强与数据去重}：将识别文本、产品字段和知识库摘要向量化，结合向量相似度匹配、规则校验和低置信字段二次补全判断是否重复入库，减少同一产品多版本记录造成的知识库污染。
        \item \textbf{服务封装}：接入 vLLM 优化模型推理服务，并使用 FastAPI 封装识别、补全、去重和入库接口，统一对外提供可复用能力，支撑商品手册、维修手册和客服知识库的持续更新。
        \item \textbf{业务落地}：构建了 1000+ 类家电产品的自动入库工作流；通过向量匹配、规则校验和二次补全的去重策略，缓解了版本迭代导致的知识冗余，降低了人工录入与重复维护成本，为下游客服问答与检索系统搭建了高质量的数据基础。
    \end{itemize}


    % \vspace{0.45em}

    %%%%%%%%%%%%%%%%%%%%
    % 奖项证书
    %%%%%%%%%%%%%%%%%%%%
    \section[在校成果]{\makebox[\iconwidth][c]{\color{primarycolor}{\faAward}}\quad 在校成果}
    % \section[证书/获奖/论文]{\makebox[\iconwidth][c]{\color{primarycolor}{\faAward}}\quad 证书/获奖/论文}
    % \section[奖项证书]{\makebox[\iconwidth][c]{\color{primarycolor}{\faTrophy}}\quad 奖项证书}
    \begin{itemize}
        % \item \textbf{证书}：CET-6
        \item \textbf{竞赛}：2025 “互联网+”大学生创新创业大赛银奖；2025 中国国际大学生创新大赛产业赛道校级二等奖；2025 首届重庆市 AI 大模型创新应用大赛省级三等奖
        \item \textbf{论文}：《LHAF: A Lightweight Hierarchical Adaptive Framework for LLM-based Multimodal Emotion Recognition in Conversations》（SCI I 区在投）
    \end{itemize}

    \vspace{0.35em}

    %%%%%%%%%%%%%%%%%%%%
    % 校园与科研经历
    %%%%%%%%%%%%%%%%%%%%

    % \section[校园与科研经历]{\makebox[\iconwidth][c]{\color{primarycolor}{\faFlask}}\quad 校园与科研经历}
    % \begin{itemize}
    %     \item \textbf{AI 竞赛实践}：参与 AI 大模型创新应用、互联网+、中国国际大学生创新大赛等竞赛项目，围绕大模型应用场景、技术方案设计与项目成果展示沉淀材料，相关项目获省级三等奖、银奖和校级二等奖。
    %     \item \textbf{校企项目实践}：在医药供应链数据问答、家电售后知识库建设等校企/实习项目中，参与需求拆解、数据处理、RAG/Prompt 链路设计、接口联调与结果评估，积累从技术方案到项目交付的完整实践经验。
    %     \item \textbf{科研训练}：围绕 LLM-based Multimodal Emotion Recognition in Conversations 方向开展科研训练，参与轻量化层次自适应框架相关论文工作，论文《LHAF: A Lightweight Hierarchical Adaptive Framework for LLM-based Multimodal Emotion Recognition in Conversations》SCI I 区在投。
    % \end{itemize}
    % \vspace{0.35em}

    %%%%%%%%%%%%%%%%%%%%
    % 自我评价
    %%%%%%%%%%%%%%%%%%%%

    % \section[自我评价]{\makebox[\iconwidth][c]{\color{primarycolor}{\faUserCheck}}\quad 自我评价}
    % \begin{itemize}
    %     \item \textbf{AI 应用落地能力}：具备 LLM、RAG、Text-to-SQL、多模态识别等项目实践，能围绕业务问题完成需求理解、方案拆解、Prompt/RAG 流程设计、服务封装与部署联调。
    %     \item \textbf{问题闭环意识}：关注上下文丢失、槽位遗漏、检索不准、SQL 生成失败、重复入库等真实工程问题，能够通过 Golden Dataset、RAGAS 指标、Bad Case 回流和 SQL 预执行等方式持续优化效果。
    %     \item \textbf{工程协作与学习迭代}：具备 Python、FastAPI、Docker、Git 等工程实践基础，熟悉接口联调、测试验证与项目复盘流程，能够结合 AI 编程工具提升开发效率并保持代码可维护性。
    % \end{itemize}

\end{document}
